\scanpage{199}
par $\mathbb{Q}_{\ell}(\unsure{k|K})$ --- c’est particulièrement intéressant dans le cas $k$ un corps de nbs.

(iv) Si $x$ est un pt à valeurs \uline{dans $\mathbb{C}$}, il lui correspond

\[
\unsure{\xi|\zeta}_{x,\mathbb{Z}}\in\mathbb{T}_{\unsure{\alpha|x}}(\mathbb{Z})\qquad
\xi_{\alpha}(\mathbb{Z})(F)=\text{réseau entier des }F_x,
\]

\deleted{\unsure{fonct|Quand} --- De plus, le groupe $\mathbb{Z}/2\mathbb{Z}$ opère sur $\mathbb{C}$ \inserted{$\xi_{x,\mathbb{Z}}$ et $\xi_{\unsure{\tau x|x},\mathbb{Z}}$ sont \unsure{can} isom.} \unsure{foncteurs par $\mathbb{Z}$} ($\tau$ est la conj. complexe),}

et des isom can

\[
\left\{\begin{aligned}
\xi_{x,\mathbb{Z}}\otimes_{\mathbb{Z}}\mathbb{Z}_{\ell}&\simeq\xi_{x,\ell}\\
\xi_{x,\mathbb{Z}}\otimes_{\mathbb{Z}}\mathbb{C}&\simeq\xi_{x,\infty}
\end{aligned}\right.
\]

D’ailleurs $\xi_{x,\mathbb{Z}}$ est défini en termes d’un pt à val. dans un corps \inserted{top} $K$ isom à $\mathbb{C}$ (à isom \unsure{can} près) en d’autres termes on a un isom canonique ``involutif''

\[
\xi_{x,\mathbb{Z}}\simeq\xi_{\tau x,\mathbb{Z}}
\]

où $\tau x$ est le pt complexe conjugué.

\deleted{Donc} On peut dire aussi, si $x$ est localisé en $a\in X$, que \inserted{a) si $k(a)$ est dense dans $\mathbb{C}$, alors} $\xi_{x,\mathbb{Z}}$ dépend de la donnée d’une topologie \inserted{archimédienne} (ou valuation) sur $k(a)$, à complété $\mathbb{C}$ \enspace b) si $k(a)$ \unsure{non dense dans|non dense ds}.

% archive sheet number: 198
